\documentclass[../cours_1re.tex]{subfiles}
\begin{document}
\setcounter{chapter}{15}
\chapter{Produit scalaire et orthogonalité}
\label{ch:16}
\begin{objectif}
Définir le produit scalaire par projection ; démontrer ses expressions et ses règles.
Prérequis : vecteurs et coordonnées de seconde, Pythagore, cercle trigonométrique.
\end{objectif}
\section{Une mesure signée de l'alignement}
Soient $\vec u=\overrightarrow{OA}$ et $\vec v=\overrightarrow{OB}$.
Si $\vec u\ne\vec0$, on projette orthogonalement $B$ sur la droite $(OA)$ en $H$.
La longueur algébrique $\overline{OH}$ est positive dans le sens de $\vec u$.
\begin{definition}
$\vec u\cdot\vec v=\|\vec u\|\overline{OH}$.
Si $\vec u=\vec0$, on pose $\vec u\cdot\vec v=0$.
Comme $\overline{OH}=\|\vec v\|\cos\theta$ pour l'angle non orienté entre
les deux vecteurs non nuls, $\vec u\cdot\vec v=\|\vec u\|\|\vec v\|\cos\theta$.
\end{definition}
Ce résultat ne dépend pas du point $O$ choisi, mais des vecteurs.
Il est positif pour un angle aigu, négatif pour un angle obtus, nul pour
un angle droit. On convient que le vecteur nul est orthogonal à tout vecteur.

\begin{center}
\begin{tikzpicture}[scale=.85,>=Latex]
\coordinate(O)at(0,0);\coordinate(A)at(5,0);\coordinate(B)at(3,2.2);\coordinate(H)at(3,0);
\draw[->,thick,accent](O)--(A)node[right]{$A$};
\draw[->,thick](O)--(B)node[above]{$B$};
\draw[dashed](B)--(H)node[below]{$H$};
\draw(3,.25)--(3.25,.25)--(3.25,0);
\node[below left]at(O){$O$};
\draw[->](.8,0)arc(0:36.25:.8);
\node at(1.15,.36){$\theta$};
\end{tikzpicture}
\end{center}

\section{La formule en coordonnées, sans circularité}
\begin{propriete}
Dans un repère orthonormé, si $\vec u=(x,y)$ et $\vec v=(x',y')$,
\[
\vec u\cdot\vec v=xx'+yy',\qquad
\|\vec u\|=\sqrt{x^2+y^2}.
\]
\end{propriete}
\begin{demonstration}
Notons $a=OA$, $p=\overline{OH}$ et $h=BH$. Pythagore donne
$OB^2=p^2+h^2$ et $AB^2=(a-p)^2+h^2$, même si $H$ est à
l'extérieur du segment $[OA]$. Par soustraction,
$AB^2=OA^2+OB^2-2ap$. Par la définition du produit scalaire,
\[
\vec u\cdot\vec v=\frac{\|\vec u\|^2+\|\vec v\|^2-\|\vec u-\vec v\|^2}{2}.
\]
Si $\vec u$ est nul, cette égalité reste vraie directement.
Dans le repère initial, Pythagore donne les trois normes au carré ;
leur développement donne $xx'+yy'$. Aucune bilinéarité n'a été supposée.
\end{demonstration}
La formule n'est pas valable sans modification dans un repère dont les axes
ne sont pas perpendiculaires ou dont les vecteurs de base ne sont pas unitaires.
Dans une base quelconque $(\vec e_1,\vec e_2)$, il faut développer
$(x\vec e_1+y\vec e_2)\cdot(x'\vec e_1+y'\vec e_2)$ en tenant compte
des produits $\vec e_i\cdot\vec e_j$.

\section{Les règles de calcul deviennent des théorèmes}
\begin{propriete}
Le produit scalaire est symétrique et bilinéaire :
\[
\vec u\cdot\vec v=\vec v\cdot\vec u,\quad
\vec u\cdot(\vec v+\vec w)=\vec u\cdot\vec v+\vec u\cdot\vec w,\quad
\vec u\cdot(\lambda\vec v)=\lambda(\vec u\cdot\vec v).
\]
De plus $\vec u\cdot\vec u=\|\vec u\|^2$, et
$\vec u\perp\vec v\iff\vec u\cdot\vec v=0$.
\end{propriete}
\begin{demonstration}
La commutativité et la distributivité des réels appliquées à $xx'+yy'$
donnent les trois premières règles. La formule de la norme donne la quatrième.
Pour deux vecteurs non nuls, le produit s'annule exactement quand
$\cos\theta=0$, donc quand $\theta=\pi/2$ ; le cas nul suit de la convention.
\end{demonstration}
Ainsi
\[
\|\vec u+\vec v\|^2=\|\vec u\|^2+2\vec u\cdot\vec v+\|\vec v\|^2,
\qquad
\|\vec u-\vec v\|^2=\|\vec u\|^2-2\vec u\cdot\vec v+\|\vec v\|^2.
\]
On les prouve en développant $(\vec u\pm\vec v)\cdot(\vec u\pm\vec v)$.
Dans une base orthonormée $(\vec i,\vec j)$, si $\vec u=x\vec i+y\vec j$,
alors $\vec u\cdot\vec i=x$ et $\vec u\cdot\vec j=y$.
Les coordonnées sont donc elles-mêmes des projections.

\begin{exemple}
Pour $\vec u=(2,1)$ et $\vec v=(-1,3)$, le produit vaut $1$, les normes
valent $\sqrt5$ et $\sqrt{10}$. L'angle est aigu et
$\cos\theta=1/\sqrt{50}$. En revanche $(1,-2)$ est orthogonal à $\vec u$.
\end{exemple}

\section{Un premier résultat d'analyse : Cauchy--Schwarz}
\approfondissement
Pour tous vecteurs $\vec u,\vec v$,
$|\vec u\cdot\vec v|\le\|\vec u\|\|\vec v\|$.
Le cosinus donne déjà une preuve dans le plan. Voici une preuve qui
fonctionnera dans des espaces beaucoup plus généraux.
\begin{demonstration}
Si $\vec v=\vec0$, l'égalité est immédiate. Sinon, pour tout réel $t$,
\[
0\le\|\vec u-t\vec v\|^2
=\|\vec u\|^2-2t(\vec u\cdot\vec v)+t^2\|\vec v\|^2.
\]
En prenant $t=(\vec u\cdot\vec v)/\|\vec v\|^2$, on trouve
$(\vec u\cdot\vec v)^2\le\|\vec u\|^2\|\vec v\|^2$.
Il y a égalité exactement lorsque $\vec u=t\vec v$, donc lorsque
les vecteurs sont colinéaires, y compris le cas nul.
\end{demonstration}
En développant $\|\vec u+\vec v\|^2$ puis en utilisant cette inégalité,
on obtient $\|\vec u+\vec v\|\le\|\vec u\|+\|\vec v\|$ :
c'est l'inégalité triangulaire.

\begin{finchapitre}
\section{Utilité et prolongements}
Le produit scalaire remplace de nombreux calculs d'angles par des calculs
algébriques. En mécanique, un travail s'écrit comme un produit scalaire
force--déplacement lorsque la force est constante.

Les chapitres 27 et 28 l'étendent aux listes de nombres et aux variables aléatoires.
Les projections deviendront alors des approximations optimales, puis, en master,
des outils d'analyse fonctionnelle et de statistique.
\begin{retenir}
Choisir entre projection, cosinus et coordonnées. Une base orthonormée
est une hypothèse essentielle. Développer une norme au carré avec le terme croisé.
\end{retenir}
\end{finchapitre}
\end{document}
